\section*{Hoofdstuk \ref{ch:g11:minerals} --- Mineralen}

\begin{solution}{exo:g11:minerals:1}
\omterm{def:g11:minerals:cleavage}{Splijting} is de neiging van een mineraal om te breken langs platte vlakken, evenwijdig aan
richtingen waarin de bindingen tussen zijn atomen zwakker zijn; breuk is
breken langs onregelmatige oppervlakken, in een mineraal zonder zwaktevlak.
Kwarts heeft geen \omterm{def:g11:minerals:cleavage}{splijting}: het vertoont een schelpvormige breuk, met gladde gebogen
oppervlakken zoals gebroken glas. Een mica heeft één \omterm{def:g11:minerals:cleavage}{splijtrichting}, perfect: het laat zich
in blaadjes afpellen.
\end{solution}

\begin{solution}{exo:g11:minerals:2}
Een spoor van een ander element is genoeg om de kleur van kwarts te veranderen, dat
helder, wit, zwart, roze, blauw of paars kan zijn; olivijn is bijna altijd
groen. Betrouwbaarder: de \omterm{def:g11:minerals:cleavage}{splijting} (of de breuk), de hardheid, de \omterm{def:g11:minerals:lustre}{glans},
de \omterm{def:g11:minerals:habit}{habitus} van vrije kristallen, de dichtheid.
\end{solution}

\begin{solution}{exo:g11:minerals:3}
(1) Olivijn. (2) Biotiet. (3) Orthoklaas (een veldspaat: twee \omterm{def:g11:minerals:cleavage}{splijtrichtingen}, hardheid
6). (4) Kwarts (geen \omterm{def:g11:minerals:cleavage}{splijting}, hardheid 7).
\end{solution}

\begin{solution}{exo:g11:minerals:4}
(1) Vijf atomen: één silicium in het midden en vier zuurstofatomen op de hoekpunten van
een tetraëder. (2) $+4 + 4 \times (-2) = -4$: de groep is \ce{SiO4^4-}.
\end{solution}

\begin{solution}{exo:g11:minerals:5}
Olivijn: losse tetraëders. Pyroxenen: enkele ketens. Amfibolen: dubbele
ketens. Mica's: lagen. Veldspaten en kwarts: raamwerken.
\end{solution}

\begin{solution}{exo:g11:minerals:6}
(1) $M(\ce{SiO2}) = 28.085 + 2 \times 15.999 = \qty{60.083}{g/mol}$;
zuurstof: $31.998/60.083 \approx \qty{53.3}{\percent}$.
(2) $M(\ce{Mg2SiO4}) = 2 \times 24.305 + 28.085 + 4 \times 15.999 =
\qty{140.691}{g/mol}$; zuurstof: $63.996/140.691 \approx \qty{45.5}{\percent}$.
(3) De korst, met \qty{46}{\percent}, is van dezelfde orde: ze bestaat bijna
helemaal uit zuurstofrijke mineralen, de silicaten.
\end{solution}

\begin{solution}{exo:g11:minerals:7}
De witte korrel heeft twee \omterm{def:g11:minerals:cleavage}{splijtrichtingen} en een hardheid tussen 5.5 en 6.5: een
veldspaat (6). De grijze korrel is harder dan 6.5 en vertoont een gebogen,
schelpvormige breuk zonder \omterm{def:g11:minerals:cleavage}{splijting}: kwarts (7).
\end{solution}

\begin{solution}{exo:g11:minerals:8}
(1) $4 - 2/2 = 3$ zuurstofatomen per silicium: \ce{SiO3}. (2) \ce{(Si,Al)2O6} bevat
twee silicium- (of aluminium)atomen voor zes zuurstofatomen, drie per silicium: het is
twee keer \ce{SiO3}.
\end{solution}

\begin{solution}{exo:g11:minerals:9}
$\rho = 0.55 \times 2.69 + 0.35 \times 3.4 + 0.10 \times 3.32 = 1.480 + 1.190
+ 0.332 \approx \qty{3.00}{g/cm^3}$. De kubus heeft een volume van
\qty{1000}{cm^3}, dus een massa van ongeveer \qty{3.00}{kg}.
\end{solution}

\begin{solution}{exo:g11:minerals:10}
(1) $0.75 \times 3.32 + 0.25 \times 3.4 = 2.49 + 0.85 = \qty{3.34}{g/cm^3}$.
(2) $(3.34 - 3.00)/3.00 \approx 0.11$: ongeveer \qty{11}{\percent} dichter.
(3) Het basalt en het \omterm{def:g11:ocean-ridges:gabbro}{gabbro} van de \omterm{def:g11:ocean-ridges:oceanic-lithosphere}{oceanische korst} zijn lichter dan de
\omterm{def:g11:ocean-ridges:peridotite}{peridotiet} van de mantel: ze blijven erbovenop liggen.
\end{solution}

\begin{solution}{exo:g11:minerals:11}
(1) Plagioklaas $168/400 = 0.42$, pyroxeen $0.33$, olivijn $0.15$, gaten
$0.10$. (2) De propositie geldt voor een gesteente zonder poriën; hier is een tiende van
het volume leeg. (3) In het vaste deel zijn de aandelen $0.42/0.90
\approx 0.467$, $0.367$ en $0.167$:
$\rho_{\text{vast}} \approx 0.467 \times 2.69 + 0.367 \times 3.4 + 0.167
\times 3.32 \approx \qty{3.06}{g/cm^3}$. Met zijn gaten, die volume toevoegen
en geen massa, heeft het gesteente de dichtheid $0.42 \times 2.69 + 0.33 \times 3.4 +
0.15 \times 3.32 \approx \qty{2.75}{g/cm^3}$, dat wil zeggen $0.9 \times 3.06$.
\end{solution}

\begin{solution}{exo:g11:minerals:12}
(1) Zuurstof $464/16.0 = 29.0$~mol; silicium $283/28.1 \approx 10.1$; aluminium
$84/27.0 \approx 3.1$; ijzer $52/55.8 \approx 0.93$; calcium $46/40.1 \approx
1.15$; magnesium $28/24.3 \approx 1.15$; natrium $23/23.0 = 1.0$; kalium
$15/39.1 \approx 0.38$. (2) In totaal ongeveer $46.8$~mol atomen, waarvan
$29.0/46.8 \approx 0.62$: bijna twee op de drie atomen zijn zuurstof.
(3) Grote zuurstofionen vormen het grootste deel van de atomen, en een nog groter deel van het
volume: de korst is, op de schaal van atomen, een stapeling van zuurstofionen, met
kleine silicium- en metaalionen in de ruimten ertussen.
\end{solution}

\begin{solution}{exo:g11:minerals:13}
(1) Het gemiddelde aantal gedeelde hoekpunten is $(2 + 3)/2 = 2.5$, dus er zijn
$4 - 2.5/2 = 2.75$ zuurstofatomen per silicium. (2) Vier siliciumatomen voor elf zuurstofatomen:
\ce{Si4O11}. (3) \ce{Si7AlO22} bevat acht tetraëders (zeven silicium, één
aluminium) voor 22 zuurstofatomen, dat wil zeggen $22/8 = 2.75$ per tetraëder, of
twee keer \ce{Si4O11}. De twee \ce{OH}-groepen horen niet bij de tetraëders.
\end{solution}

\begin{solution}{exo:g11:minerals:14}
$0.20 \times 2.62 + 0.10 \times 2.56 + 0.50 \times 2.65 + 0.10 \times 3.09 +
0.10 \times 3.23 = 0.524 + 0.256 + 1.325 + 0.309 + 0.323 \approx
\qty{2.74}{g/cm^3}$, boven \qty{2.7}{g/cm^3}; het graniet van het voorbeeld,
rijker aan kwarts en kaliveldspaat, is ongeveer \qty{2.59}{g/cm^3}, onder
\qty{2.6}{g/cm^3}. De gesteenten van het westen zijn dichter omdat ze meer
plagioklaas en meer biotiet en hoornblende bevatten, rijk aan ijzer en magnesium; de
silicarijke granieten bevatten meer kwarts en kaliveldspaat, de lichtste
gesteentevormende mineralen.
\end{solution}

\begin{solution}{exo:g11:minerals:15}
(1) Bovenkorst: $66.6 \times 0.467 \approx \qty{31.1}{\percent}$ silicium;
hele korst: $60.6 \times 0.467 \approx \qty{28.3}{\percent}$. (2) De bovenkorst,
rijker aan silica, zou meer kwarts en veldspaten moeten bevatten, en minder
mineralen die rijk zijn aan ijzer en magnesium; de diepere korst het omgekeerde. (3) Kwarts
en veldspaten zijn de lichtste gesteentevormende mineralen: de bovenkorst zou
lichter moeten zijn dan de diepere korst.
\end{solution}

\begin{solution}{pb:g11:minerals:1}
\textbf{1.} Kwarts: geen \omterm{def:g11:minerals:cleavage}{splijting}, schelpvormige breuk, glasglans, harder
dan de spijker (7).

\textbf{2.} Orthoklaas, een kaliveldspaat: twee \omterm{def:g11:minerals:cleavage}{splijtrichtingen}, hardheid 6 op
de schaal van Mohs.

\textbf{3.} Plagioklaas (hier oligoklaas), de veldspaat met natrium en
calcium.

\textbf{4.} D is biotiet, E muscoviet: twee mica's. Hun tetraëders vormen
lagen, en de bindingen tussen de lagen zijn zwak: ze splijten langs deze ene
perfecte \omterm{def:g11:minerals:cleavage}{splijtrichting}.

\textbf{5.} Kwarts en de twee veldspaten (A, B, C) zijn raamwerken; biotiet
(D) bevat ijzer en magnesium.

\textbf{6.} Kwarts $150/500 = 0.30$; orthoklaas $0.35$; oligoklaas $0.25$;
biotiet $0.06$; muscoviet $0.04$.

\textbf{7.} Wanneer de korrels zonder orde verspreid liggen, schat het aandeel punten
op een mineraal het aandeel van de oppervlakte dat het op het vlak bedekt, en dat
is gelijk aan het aandeel van het volume dat het vult.

\textbf{8.} $0.35 + 0.25 = 0.60$: \qty{60}{\percent} veldspaat.

\textbf{9.} Met $1/500 = 0.002$, dat wil zeggen \qty{0.2}{\percent}.

\textbf{10.} Ja: veldspaat, kwarts en mica.

\textbf{11.} $\rho = \phi_1 \rho_1 + \phi_2 \rho_2 + \dots + \phi_k \rho_k$,
met $\phi_1 + \dots + \phi_k = 1$.

\textbf{12.} Kwarts $0.30 \times 2.62 = 0.786$; orthoklaas $0.35 \times 2.56 =
0.896$; oligoklaas $0.25 \times 2.65 = 0.6625$; biotiet $0.06 \times 3.09
\approx 0.185$; muscoviet $0.04 \times 2.82 \approx 0.113$ (alle in
\unit{g/cm^3}).

\textbf{13.} $0.786 + 0.896 + 0.6625 + 0.185 + 0.113 \approx
\qty{2.64}{g/cm^3}$.

\textbf{14.} Orthoklaas, met $0.896$, omdat het het meest voorkomt. Biotiet
en muscoviet, met \qty{3.09}{} en \qty{2.82}{g/cm^3}, zijn dichter dan de
andere mineralen en verhogen de dichtheid van het gesteente tot boven die van kwarts en
veldspaat alleen.

\textbf{15.} \qty{2640}{kg/m^3}: \qty{1}{m^3} van het graniet heeft een massa van
ongeveer \qty{2.64}{t}.

\textbf{16.} $100 \times 2.643 \approx \qty{264.3}{g}$.

\textbf{17.} Het blok verplaatst \qty{100}{cm^3} water, met een massa van
\qty{100}{g}: de balans geeft $264.3 - 100 = \qty{164.3}{g}$ aan.

\textbf{18.} $264.3/(264.3 - 164.3) = 264.3/100 \approx \qty{2.64}{g/cm^3}$.

\textbf{19.} Tussen de twee: lichter dan het gemiddelde van \qty{2.69}{g/cm^3},
boven de \qty{2.6}{g/cm^3} van de granieten die het rijkst zijn aan silica en kalium.
Meer plagioklaas en meer biotiet of hoornblende, en minder kwarts en
kaliveldspaat, zouden een graniet dichter maken.

\textbf{20.} De dichtheid van het graniet is ongeveer
\textbf{\qty{2.64}{g/cm^3}}.
\end{solution}
